% ===================================================================
%  TABEL Nr. 10 — Meri. Sadam.
%  Traduction 2026 d'apres texto/io, controlee sur texto/fr et
%  texto/fi. Consigne : voir texto/et/10-tabel-01.tex.
%
%  LE TABLEAU 10 ROMANCHE ETAIT ECRIT PAR UNE LANGUE SANS MER ; CELUI-
%  CI EST ECRIT PAR UNE LANGUE QUI EN A TROIS MILLE HUIT CENTS
%  KILOMETRES. ET LE RESULTAT N'EST PAS CELUI QU'ON ATTENDAIT.
%
%  Cette page porte trois ages d'emprunt, et on les distingue a
%  l'oeil :
%
%  1. L'EMPRUNT INVISIBLE. « laev » le navire et « paat » la barque
%     sont germaniques, entres il y a plus de mille ans. Ils se
%     declinent comme n'importe quel mot du fonds, ils font des
%     composes — « laevatehas » le chantier, « paadimees » le
%     batelier — et rien dans leur forme ne les denonce. Un locuteur
%     ne les sent pas plus etrangers que « meri » ou « rand ».
%
%  2. L'EMPRUNT VISIBLE. « kruiiser », « torpeedo », « fregatt »,
%     « admiral », « arsenal », « semafor ». Ils sont entres au XIXe
%     siecle, ils gardent leur forme d'origine et leurs voyelles
%     longues doublees les signalent a chaque ligne.
%
%  3. CE QUI N'A PAS ETE EMPRUNTE DU TOUT : « meri », « rand »,
%     « laine », « tuul », « saar », « aer », « võrk », « kalur ».
%     La mer, la rive, la vague, le vent, l'ile, la rame, le filet,
%     le pecheur.
%
%  LA TROISIEME LISTE EST CELLE DE L'HOMME QUI POSSEDE SON BATEAU ; LA
%  DEUXIEME EST CELLE DE L'HOMME QUI SERT SUR CELUI D'UN AUTRE. Le
%  tableau 7 l'avait montre sur le moulin et le meunier ; ici c'est la
%  meme frontiere, tracee sur une cote entiere.
%
%  MAIS LA PREMIERE LISTE AJOUTE CE QUE LE MOULIN NE POUVAIT PAS
%  DONNER : L'AGE D'UN EMPRUNT SE LIT A CE QU'IL A CEDE. Un mot pris
%  tot se plie a la morphologie de la langue et disparait dans le
%  fonds ; un mot pris tard reste entier et se voit. La colonne
%  romanche avait montre au tableau 2 qu'on ne peut pas lire LA ROUTE
%  d'un emprunt dans sa forme ; l'estonien montre qu'on peut y lire
%  SA DATE.
% ===================================================================

%%K t10-apar-1 apar
\VUsaut{4.0mm}
\VUtitre{15.0pt}{60}{TABEL Nr. 10}
\VUsaut{3.0mm}
\VUpk{11.4pt}{40}{Meri. --- Sadam.}
\VUsaut{3.0mm}

%%K t10-01-1 p
1. --- Suvel, sellal kui ühed otsivad mäel rahu ja jahedust, eelistavad teised minna
rannikule. Nii tegime meie eelmisel aastal, sest me armastame väga \VUgras{ookeani}
\textsuperscript{(1)}.

%%K t10-02-1 p
2. --- Mina tunnen suurt lõbu, vaadeldes seda tohutut veeavarust, mis ulatub
\VUgras{silmapiirini} \textsuperscript{(2)}, ja nähes pikale \VUgras{ülesõidule}
väljumas hiiglaslikke \VUgras{aurulaevu} \textsuperscript{(3)}, millele lähevad pardale
Ameerikasse sõitvad reisijad.

%%K t10-03-1 p
3. --- Seepärast valib mu isa, kes tunneb mu kalduvust, vaheaja veetmiseks alati mõne
väga rahuliku ranna, mis asub ühe meie suure \VUgras{sõjasadama} lähedal.

%%K t10-03-2 p
Meie viimase peatuse ajal tuli \VUgras{eskaader} ankrusse tohutu \VUgras{muuli}
\textsuperscript{(4)} taha, mis sulgeb ja kaitseb \VUgras{sõjasadamat}
\textsuperscript{(5)}, ja mul õnnestus oma soovi järgi imetleda mitmesuguseid
\VUgras{sõjalaevu}, väikest \VUgras{allveelaeva} \textsuperscript{(10)}, mis sukeldub ja
kaob vete alla, ja ka tohutut \VUgras{soomuslaeva} \textsuperscript{(9)}, mis seni kartis
peaaegu ainult \VUgras{torpeedopaadi} \textsuperscript{(8)} rünnakuid.

%%K t10-tit-2 sub
\VUsaut{3.0mm}
\VUpk{10.2pt}{40}{Väikesed laevad.}
\VUsaut{3.0mm}

%%K t10-04-1 p
4. --- See viimane oskas juba väga \VUgras{osavalt} leida paksu metallsoomuse nõrga
koha, et sinna torgata ohtlik \VUgras{torpeedo}, mille plahvatus toob kaasa laeva
hukkumise, kuid siiski oli see vähem kardetav kui allveelaev, sest hävitajad ajavad
seda kergesti taga.

%%K t10-05-1 p
5. --- Ma võisin näha ka kiireid soomustatud \VUgras{ristlejaid} \textsuperscript{(6)},
peaaegu sama haavamatuid kui tõeline soomuslaev, mitut \VUgras{transpordilaeva}
\textsuperscript{(12)}, kolme või nelja \VUgras{kahurpaati} \textsuperscript{(7)} ja
lõpuks üht \VUgras{fregatti} \textsuperscript{(11)}. \VUgras{Selle laevastiku vaatepilt,
kui see ankru hiivamiseks manööverdab, on kõige huvitavam.}

%%K t10-06-1 p
6. --- Kui suur ehituse erinevus nende suurte teraskuhjade ja elegantsete
\VUgras{kolmemastiliste} või graatsiliste \VUgras{purjelaevade} \textsuperscript{(13)}
vahel, mida kaubalaevastikus veel praegugi kasutatakse ja mida ma nägin kõigi
purjedega sadamasse sisenemas, kui ma jalutasin \VUgras{kail} \textsuperscript{(14)}.
Tõsi küll, need kaotasid palju oma eelistest, \VUgras{kui} \VUgras{tuul} oli vastu. Nad
pidid kõik purjed alla laskma, et basseini tagasi tulla, ja \VUgras{puksiiri}
\textsuperscript{(16)} oli vaja, et neid ära tuua ja kai äärde tuua nagu lihtsaid
\VUgras{praame} \textsuperscript{(17)}.

%%K t10-tit-3 sub
\VUsaut{3.0mm}
\VUpk{10.2pt}{40}{Kaubasadam.}
\VUsaut{3.0mm}

%%K t10-07-1 p
7. --- Selles sadamas, millest ma räägin, on huvitav see, et see ei sisalda ainult
sõjasadamat, mis on hiiglaslike töödega kunstlikult rajatud, vaid ka imepärast
\VUgras{kaubasadamat}, mis asub suure jõe suudmes.

%%K t10-08-1 p
8. --- Maakeel, mis eraldab kaht basseini, on kindlustatud ja kaitstud
\VUgras{patareide} ja \VUgras{bastionide} reaga, mis ulatuvad merre. Vaja on eriluba, et
jalutada \VUgras{vallidel} \textsuperscript{(15)}. Ma sain selle kergesti oma onu kaudu,
kes on \VUgras{laevatehaste} \textsuperscript{(18)} insener, ja ma läksin vaatama laevu,
mida ehitatakse avarate angaaride all.

%%K t10-09-1 p
9. --- Ma viibisin isegi soomuslaeva vettelaskmisel ja ma mäletan liigutavat hetke,
mida kogu rahvahulk tundis, kui hiiglaslik mass, endale jäetud, hakkas libisema oma
hälli sarnasel raamil ja tuli jõe veele ujuma.

%%K t10-10-1 p
10. --- Laevatehastesse minekuks ületatakse \VUgras{õppeväli} \textsuperscript{(19)}, kus
garnisoni sõdurid tulevad iga päev harjutama, ja minnakse üle raudse \VUgras{sillakese}
\textsuperscript{(20)}, mis ühendab paraadiplatsi \VUgras{arsenaliga}
\textsuperscript{(21)}.

%%K t10-tit-4 sub
\VUsaut{3.0mm}
\VUpk{10.2pt}{40}{Arsenal ja dokid.}
\VUsaut{3.0mm}

%%K t10-11-1 p
11. --- Koos oma onuga külastasin ma seda suurt ehitist, mis on täis \VUgras{kahureid}
\textsuperscript{(22)}, \VUgras{kuule} \textsuperscript{(23)} ja \VUgras{mürske}. Ma nägin
ka \VUgras{metallitehast} \textsuperscript{(24)}, kus valmistatakse aurulaevade
\VUgras{katlaid} \textsuperscript{(25)}.

%%K t10-12-1 p
12. --- Väga lähedal on \VUgras{dokid} \textsuperscript{(26)} --- parandusdokid --- ja
väga tähelepanuväärsed \VUgras{ujuvdokid} \textsuperscript{(27)}, kus ma võisin väga
lähedalt näha, parandatava laeva peal, paadi \VUgras{sõukruvi} \textsuperscript{(28)},
\VUgras{kiilu} \textsuperscript{(29)} ja \VUgras{rooli} \textsuperscript{(30)}.

%%K t10-13-1 p
13. --- Kaubasadam on täpselt sama uurimist väärt kui sõjasadam, ja ma veetsin palju
tunde kaidel ammuli sui, kus on kinnitatud \VUgras{rattaaurikud}
\textsuperscript{(31)}, mis sõidavad mööda jõge üles, ja vaadates töötamas
\VUgras{mastkraanasid} \textsuperscript{(32)}, mis tõstavad nagu sulgi tohutuid koormaid.

%%K t10-tit-5 sub
\VUsaut{4.0mm}
\VUpk{10.2pt}{40}{Loots.}
\VUsaut{3.0mm}

%%K t10-14-1 p
14. --- Ma imetlesin \VUgras{ookeaniaurikuid} \textsuperscript{(34)}, mis väljusid
reidilt, oma \VUgras{lippudega} \textsuperscript{(35)} tuules, juhituna osava
\VUgras{lootsi} \textsuperscript{(36)} poolt, kes oskab imepäraselt kulgeda mööda
faarvaatrit, mis on tähistatud \VUgras{poidega} \textsuperscript{(40)} ja märkidega,
vältides \VUgras{lodju} \textsuperscript{(41)}, mida on raske juhtida nende ainsa
nelinurkse purjega. Ma eristasin \VUgras{vööris} \textsuperscript{(37)} --- eesosas ---
teise klassi \VUgras{kajutid} \textsuperscript{(39)}, ja \VUgras{ahtris}
\textsuperscript{(38)} --- tagaosas --- esimese klassi salongid.

%%K t10-15-1 p
15. --- Mõnikord viibisin ma ka \VUgras{kaubalaeva} \textsuperscript{(42)} laadimisel,
millesse oli äsja kuhjatud \VUgras{lao} \textsuperscript{(43)} ja \VUgras{merejaama}
\textsuperscript{(44)} kaup, mille olid kohale toonud \VUgras{kaeraudtee}
\textsuperscript{(45)} arvukad rongid.

%%K t10-tit-6 sub
\VUsaut{3.0mm}
\VUpk{10.2pt}{40}{Sõudmise lõbu.}
\VUsaut{3.0mm}

%%K t10-16-1 p
16. --- Sageli sõudsin ma \VUgras{ujuvdokis} \textsuperscript{(53)}, kuhu tulevad varjuma
\VUgras{paadid} \textsuperscript{(46)}, ja mulle oli rõõm käsitseda \VUgras{aeru}
\textsuperscript{(47)}. Ma tõusin \VUgras{tekile} \textsuperscript{(48)} ühel
\VUgras{purjepaadil} \textsuperscript{(49)}, ma ronisin mööda \VUgras{trosse}
\textsuperscript{(52)} \VUgras{mastile} \textsuperscript{(51)} kuni \VUgras{raadeni}
\textsuperscript{(50)}, siis jätkasin ma oma jalutuskäiku \VUgras{joolaga}
\textsuperscript{(54)} raskete \VUgras{kalapaatide} \textsuperscript{(55)} vahel, kus
\VUgras{võrgud} \textsuperscript{(56)} kuivavad.

%%K t10-17-1 p
17. --- \VUgras{Kail} \textsuperscript{(58)} liikus mu ees kõige mitmekesisem
\VUgras{kaup} \textsuperscript{(59)}, mis oli \VUgras{kastides} \textsuperscript{(60)} või
\VUgras{pallides} \textsuperscript{(61)}. Need moodustasid lodjade \VUgras{lasti}
\textsuperscript{(63)}, ja võimsad \VUgras{kraanad} \textsuperscript{(62)} võtsid need ja
panid \VUgras{veoautodele} \textsuperscript{(64)}, sellal kui voorimehed deklareerisid
need ja maksid maksud \VUgras{tollis} \textsuperscript{(65)}.

%%K t10-18-1 p
18. --- Lõpuks, kui ma jõudsin \VUgras{pöördsillani} \textsuperscript{(66)} või
\VUgras{lüüsini} \textsuperscript{(69)}, möödudes \VUgras{süvendajast}
\textsuperscript{(68)}, mis puhastab \VUgras{kanalit} \textsuperscript{(67)}, astusin ma
maha kaile \VUgras{semafori} \textsuperscript{(70)} juures.

%%K t10-19-1 p
19. --- Selle kai teisel pool tulevad tõesti randuma \VUgras{šalupad}
\textsuperscript{(71)}, mis toovad maale eskaadri ohvitsere. On imetlusväärne vaatepilt
näha madruseid taktis oma aere tõstmas, sellal kui \VUgras{laine} \textsuperscript{(72)}
kantud paat randub kergesti maabumiskoha trepi juures. Sellel kohal saab paremini
jälgida \VUgras{loodet} ning \VUgras{tõusu} ja \VUgras{mõõna} liikumist \VUgras{rannal}
\textsuperscript{(73)}.

%%K t10-tit-7 sub
\VUsaut{3.0mm}
\VUpk{10.2pt}{40}{Klubi.}
\VUsaut{3.0mm}

%%K t10-20-1 p
20. --- Koju tagasi minekuks kasutasin ma \VUgras{elektritrammi}
\textsuperscript{(74)}, mille \VUgras{peatus} \textsuperscript{(75)} on \VUgras{posti}
juures, mis on pandud ohvitseride klubi ette.

%%K t10-20-2 p
Klubi \VUgras{rõdult} \textsuperscript{(76)}, mis on ehitud \VUgras{ankrute}
\textsuperscript{(77)}, kahurite ja kuulidega, nägin ma sageli suurepärast mereohvitseride
kogunemist: \VUgras{leitnante} \textsuperscript{(78)} oma mõõgaga, suurtükiväe
\VUgras{kapteneid} \textsuperscript{(79)} ja \VUgras{jalaväe} omi \textsuperscript{(80)} oma
\VUgras{mõõgaga}. Nende taga oli \VUgras{madrus} \textsuperscript{(81)}, kes tõi neile
\VUgras{pikksilma}, kui nad tahtsid eristada, mis kaugel toimub.

%%K t10-21-1 p
21. --- Ma nägin ka noori \VUgras{nooremleitnante} \textsuperscript{(82)}, kes said
juhiseid oma \VUgras{komandörilt} \textsuperscript{(83)} või \VUgras{admiralilt}
\textsuperscript{(84)}, sellal kui \VUgras{pootsman} \textsuperscript{(85)} ootas, et neid
laevale viia.

%%K t10-22-1 p
22. --- Sellelt rõdult on vaade suurepärane: valitsetakse kogu randa oma
\VUgras{telkide} \textsuperscript{(86)} ja \VUgras{kabiinidega} \textsuperscript{(87)}, ja
nähakse supluse ajal \VUgras{suplejaid} \textsuperscript{(88)}, kes rõõmsalt hullavad
\VUgras{kasiino} \textsuperscript{(89)} ees.

%%K t10-23-1 p
23. --- Ranna otsas moodustab rannik väikese \VUgras{kaljuse}
\VUgras{poolsaare} \textsuperscript{(91)}, kus \VUgras{murdlained} \textsuperscript{(92)}
tulevad murduma ja millele on ehitatud \VUgras{tuletorn} \textsuperscript{(93)}, mis öösel
näitab meresõitjatele teed. See poolsaar on maaga ühendatud ainult väga kitsa
\VUgras{maakitsuse} kaudu \textsuperscript{(90)}.

%%K t10-tit-8 sub
\VUsaut{3.0mm}
\VUpk{10.2pt}{40}{Rannikute üldvaade.}
\VUsaut{3.0mm}

%%K t10-24-1 p
24. --- \VUgras{Pank} \textsuperscript{(94)} laiub, meri on selle lõhestanud ja joonistab
sellesse kapriisselt rea \VUgras{lahtesid} \textsuperscript{(95)} ja \VUgras{neemi},
tehes selle väga maaliliseks.

%%K t10-24-2 p
\VUgras{Platool} \textsuperscript{(96)}, mis valitseb \VUgras{väina}
\textsuperscript{(98)}, on väga tähtis \VUgras{kindlus} \textsuperscript{(97)} ja mõned
«villad».

%%K t10-25-1 p
25. --- See väin eraldab mandrist väga ohtliku \VUgras{saare} \textsuperscript{(99)},
mida ümbritsevad \VUgras{rahud} \textsuperscript{(100)} ja \VUgras{murdlained}, kus paadid
ja isegi suured laevad mõnikord karile jooksevad. Vaatamata abi kiirusele ja
\VUgras{päästepaatide} kiirele saabumisele on need \VUgras{laevahukud} hirmuäratavad, ja
pööripäeva \VUgras{tormide} ajal on mitu laeva hukkunud koos inimeste ja kaubaga.

%%K t10-25-2 p
Vaatamata sellele naabrusele on \VUgras{sadam} madruste seas väga sagedasti külastatud.

\VUsaut{6.0mm}
\VUfilet{22mm}
